\documentclass[../main.tex]{subfiles}
\usepackage{amssymb}
\usepackage{amsmath,amsthm} 
\usepackage{amsfonts} 
\usepackage{graphicx} 
\usepackage{auto-pst-pdf} 
\graphicspath{{\subfix{../pictures/}}}
\begin{document}
\section{Závěr}
Výsledkem mé semestrální práce je program vykreslující graf zadané funk-ce který je časově nenáročný a v běžných případech je jeho běh ukončen úspě-chem za méně než 1 vteřinu. Ovšem výsledný graf není perfektní z důvodu zvoleného přístupu k vykreslení grafu, a to v ohledu na velmi rychle rostoucí funkce, kde se občas stane, že graf není dokreslen až k okrajům jako na obrázku \ref{ref_tan}.
\vspace{0.5em}

Zadání jsem splnil v plném rozsahu. Bylo by ovšem vhodné upravit zápis do souboru a validační funkce, které by bylo možné značně zjednodušit.
\vspace{0.5em}

Během řešení této úlohy jsem narazil na problémy během vykreslování funkcí. Zejména nespojité funkce a funkce definované pouze na určitém definičním oboru. Na začátku řešení jsem měl také problém s převáděním funkce do tokenů, a to protože jsem špatně použil ukazatel.

\end{document}